\section{Two divisibilities and an index quotient in one Pell root}
\label{sec:pell-carriers}

This section gives a second, direct route to smaller equations. It uses the
authentic compact native tuple and ordinary Pell kernel, and does not use
Jones's nine-variable compression. The new result has fourteen strictly
positive unknowns, exact degree $26$, and a $92$-operation evaluator. A
companion has fourteen positive unknowns, degree $30$, and $85$ operations.
The inspected wave-eight fourteen-positive point is degree $52$ at $79$
operations. Thus the first construction halves its degree and adds thirteen
gates; the second uses six additional gates and reduces its degree by
$42.3\%$. These are tradeoffs within the pinned collection, with no claim
of a global simplicity minimum.

The structural step combines the integral-root divisibility carrier of
package~58 and the main-root index absorption of package~59. A single
existing Pell root now enforces a weak congruence first, determines the
Pell index, and only then enforces a stronger congruence restoring another
divisibility. This ordered recovery permits the coordinate $X$ to be
supplied directly while deleting the old index quotient $h$.

\subsection{Fixed coefficients and exact equation}

Use the \emph{whole} package-44 compact tuple, whose binding is included
under \texttt{pell/coefficients/}. Its identifier is
\texttt{wang943-complement67-shared-anchor-length8807-prime581197}.
Set
\[
 B=2^{590069},\qquad \beta=B-1,\qquad
 K_0=D_C+BD_R,\qquad M_F=M_{F,\mathrm{native}}+B-1.
\]
The fourth port is $M_C$. The high mask $M_F$ is already shifted. The four
little-endian binary integers occupy $295{,}388$ raw bytes and $9{,}445$
gzip bytes. Their recipes, payloads, exact hashes, literal machine data,
native adapter, and independent coefficient auditor accompany the paper.
All four ports are fixed numerals, not quantified coordinates.

The fourteen positive inputs are
\[
 j,F,\alpha,z,f,i,T,x,Y,r,\eta,\zeta,u,\gamma.
\]
Define polynomial abbreviations
\begin{align*}
 J&=3j+1,& q&=\beta J+1,& C&=q-F-\alpha,&X&=x+q,\\
 E&=XY,& k&=\eta+\zeta,&c&=kY+\eta,& a&=E+Y,\\
 H&=4a+3,&\Delta&=a^2+H=(a+2)^2-1,&\ell&=EkY.
\end{align*}
Also set
\begin{align*}
 R&=\bigl[(q-1)(q-F)+\alpha\bigr](q^2-1)
                      +(M_C+qM_F)J,\\
 S&=4qY,\qquad M=qS=4q^2Y,\qquad v=k-R-1,\\
 D_*&=M(X+ac)+H(M\gamma-v+SX),\\
 U&=\Delta(f^2-1),\qquad I=ic^2,\qquad V=f^2(cT-R)-c.
\end{align*}
The five residuals are
\begin{align}
 A_0&=[q^3(\ell+r)+Y]^2-q^6[\ell(\ell+k)+1],
                                                    \label{eq:pell-first}\\
 A_m&=D_*^2-M^2(\Delta c^2+1),                       \label{eq:pell-main}\\
 A_t&=q[K_0C-F-z(q-1)]+XC,                          \label{eq:pell-transport}\\
 A_a&=(V+u)^2-Uu(2V+u)-1,                           \label{eq:pell-aux}\\
 A_s&=I^2-U.                                       \label{eq:pell-strong}
\end{align}
The single equation is
\begin{equation}
 \boxed{P_{14,26}=A_0^2+A_m^2+A_t^2+A_a^2+A_s^2=0.}
 \label{eq:pell26}
\end{equation}
There are no variable exponents, radicals, divisions or extra quantified
registers in this definition. The supplied literal circuit is
\texttt{pell/circuits/p14\_degree26.json}.

\begin{theorem}[Fourteen-positive equation]
\label{thm:pell26}
For the stated authentic coefficient tuple, $P_{14,26}$ has exact total
degree $26$ and a $92$-gate evaluator. The new reductions below, composed
with the explicitly supplied ordinary native compiler and calibration,
establish by ordinary mathematics formalizable in ZFC that
\[
 \ZFC\vdash
 \bigl[\neg\exists\mathbf t\in\Zp^{14}\;
                   P_{14,26}(\mathbf t)=0\bigr]
       \longleftrightarrow\Con(\ZFC).
\]
This is not a newly kernel-checked end-to-end formalization or a supplied
derivation file in the object-language ZFC calculus.
\end{theorem}

\subsection{The two-stage carrier lemma}

The elementary fact behind the construction is that a rational square root
of an integer is integral: $m^2\mid n^2$, with $m>0$, implies $m\mid n$.
At a zero of $A_0$, it gives
\[
 q^3\mid Y,\qquad s=Y/q^3>0,\qquad
 \tau=\ell+r+s>0,\qquad \tau^2=\ell(\ell+k)+1.
\]
Since $590069$ is odd, $\beta\equiv1\pmod3$. Thus the supplied
$J=3j+1$ gives $q\equiv2\pmod3$. As $q\mid Y$,
\[
 H=4Y(X+1)+3\equiv3\pmod q,
 \qquad \gcd(q,H)=1.
\]
The odd integer $H$ also satisfies $\gcd(Y,H)\mid3$. Hence
\begin{equation}
 \gcd(S,H)\mid3,\qquad
 \mu:=\frac S{\gcd(S,H)}\ge\frac{4qY}{3}.
 \label{eq:pell-mu}
\end{equation}
For $Y\ge q^3$ and $R<q^4-2$, this yields $\mu>R+1$.
Moreover both $S$ and $M$ divide $4XY^2$; for $M$ the quotient is
$XY/q^2=Xsq$, which is already integral without assuming $q\mid X$.

At a zero of $A_m$, the same rational-root lemma gives the integer
$g=D_*/M$. Reducing $Mg=D_*$ only modulo $S$ gives
\begin{equation}
 S\mid H(k-R-1),\qquad \mu\mid k-R-1.
 \label{eq:pell-weak}
\end{equation}
The Pell argument below uses this weaker congruence to recover the exact
index and show $M\mid v$. Only then do we return to the full modulus:
\begin{equation}
 M\mid HSX\quad\Longrightarrow\quad q\mid HX
       \quad\Longrightarrow\quad q\mid X.
 \label{eq:pell-strongdiv}
\end{equation}
The order of these steps prevents a circular use of the quotient $X/q$.

\subsection{Soundness: recovering the indices before the quotients}

At a zero of \eqref{eq:pell26}, all five residuals vanish. Restore $s$ and
$\tau$ as above. The cleared transport equation becomes the rational
equality
\[
 (K_0+X/q)C=F+z(q-1).
\]
Both its right side and its coefficient of $C$ are positive. Therefore
$C>0$ before $X/q$ has been shown integral. The authentic mask bounds give
\begin{gather}
 q\ge16,\qquad q(q^2-1)<R<q^4-q^3+q^2-1<q^4-2,\label{eq:pell-bounds}\\
 R-2>3q+1,\qquad E>q^4>R+2,\qquad a>R+2,
 \qquad a+2<2E.\nonumber
\end{gather}
For clarity, the bounds on $R$ follow from
\[
 G=(q-1)(q-F)+\alpha=q\alpha+(q-1)C,
 \quad q\le G\le q^2-q,
\]
and $0<(M_C+qM_F)J<(q-1)(1+2q)$. The supplied shift $X=x+q>q$
is what proves $E>q^4$ at this stage. The inequalities
$kY<c<k(Y+1)$ follow directly from $\eta,\zeta>0$.

The integer $g=D_*/M$ is positive. Indeed, $v<k$,
$b_0=X+ac>akY$ and $H<5a$ give
\[
 D_*>Mb_0-Hk>ak(MY-5)>0.
\]
Hence $g^2-\Delta c^2=1$. The first and main norms give positive Pell
indices $n,p$ with
\begin{align*}
 \tau&=\chi_{P}(n),&k&=2\psi_{P}(n),&P&=2XY^2+1,\\
 g&=\chi_A(p),&c&=\psi_A(p),&A&=a+2.
\end{align*}
Here $\chi_A(t)+\psi_A(t)\sqrt{A^2-1}
=(A+\sqrt{A^2-1})^t$. The first norm has fundamental integer-coordinate
solution $(2XY^2+1,2)$ for discriminant $XY^2(XY^2+1)$: a second
coordinate $1$ is excluded by consecutive squares. Since $P>A$ and $c>k$,
monotonicity implies $p>n$.

The Pell recurrence gives $k\equiv2n\pmod{4XY^2}$. Since $S\mid4XY^2$,
\eqref{eq:pell-weak} implies
\[
 2n=R+1+t\mu\quad(t\in\mathbb Z).
\]
The inequalities $0<R+1<\mu$ and $n>0$ force $t\ge0$. Thus $n>24$ and
$p\ge26$. Consequently
\begin{equation}
 c>A\Delta^2,\qquad c>2p,\qquad c>2R.
 \label{eq:pell-rankbounds}
\end{equation}
For the first estimate it suffices that $p\ge6$ and
$\psi_A(6)=32A^5-32A^3+6A>A(A^2-1)^2$. The second is elementary Pell
growth, and the third follows from $c\ge\psi_A(2)=2A>2R$.

The equation $A_s=0$ gives $U=I^2$ with $I=ic^2>1$, hence $f\ge2$.
The \emph{ordinary strong-rank lemma} supplied in
\texttt{pell/upstream/ORDINARY\_KERNEL.md}, Sections 4--5, now applies
with all its hypotheses established. It gives an index $m$ such that
\begin{equation}
 f=\chi_A(m),\quad p\mid m,\quad c\mid m,\quad m>2p,
 \quad I=\Delta\psi_A(m),\quad c^2\mid f^2-1.
 \label{eq:pell-strongrank}
\end{equation}
The last divisibility is a conclusion of the rank lemma; it cannot be
obtained by informally cancelling $\Delta$ in $I^2=\Delta(f^2-1)$.
The copied proof handles the possibly nonfundamental squarefree Pell
parameter and is an explicit inherited dependency.

The auxiliary equation is exactly
\[
 I^2V^2-(I^2-1)(V+u)^2=1.
\]
It excludes $V=0$ and $V+u=0$. Put $y=|V+u|>0$. The supplied quotient
and \eqref{eq:pell-strongrank} imply
\[
 V\equiv-c\pmod f,\qquad V\equiv-R\pmod c.
\]
Also $c^2\mid I$. The signed auxiliary-index lemma in the same ordinary
source gives an auxiliary index congruent to $\pm p$ and $\pm R$ modulo
$c$. The bounds $0<p,R<c/2$ force $p=R$. This argument does not assume
$V>0$. Its ordinary source quotient is restored by $T_{\rm old}=fT>0$.

Now $n<p=R$, so $|2n-R-1|<R<\mu$. Together with the congruence modulo
$\mu$ this yields $2n=R+1$. Therefore
\[
 v=k-R-1=k-2n
\]
is a positive multiple of $4XY^2$. Positivity follows from $n\ge2$ and
$\psi_P(n)>n$. In particular $M\mid v$, and
\eqref{eq:pell-strongdiv} proves $q\mid X$. Restore
\begin{equation}
 w=X/q>0,\qquad h=v/E>0,\qquad
 \gamma_{\rm old}=\gamma-v/M+w.
 \label{eq:pell-restore}
\end{equation}
These are integers. The main root becomes
$g=X+ac+H\gamma_{\rm old}$. Since $n\ge2$ gives $k\ge4P>X$, one has
$c>X$; and $g>\sqrt\Delta c>(a+1)c$. Thus
$\gamma_{\rm old}>(c-X)/H>0$. The transport equation is now the genuine
integer equality $(K_0+w)C=F+z(q-1)$, corresponding to old source
coordinate $z_{\rm old}=z+1$.

All positive ordinary source coordinates are recovered and all six source
factors equal $1$. The supplied native compiler and calibrated machine
endpoint therefore yield a finite refutation of the specified ZFC calculus.

\subsection{Completeness and the positive replacement witnesses}

Conversely, a finite refutation gives canonical normalized witnesses from
the authentic compact adapter. Transfer the strong coordinate by
$i=\Delta i_{\rm normalized}$, as in package~58. Then
$I^2=\Delta(f^2-1)$, $c\mid I$, and $f^2\equiv1\pmod c$.
Canonical indices satisfy
\[
 p=R=2n-1,\qquad n\ge3,\qquad R\equiv3\pmod4,
 \qquad X=2^R,\qquad Y=sq^3,
\]
and $k>XY$. These are inherited compiler/Pell facts, not conclusions of a
finite numerical search.

The adapter supplies $q=B^N$ and $J=1+B+\cdots+B^{N-1}$ with odd
$N=581197^e>1$. Since $B\equiv-1\pmod3$, $J\equiv1\pmod3$ and
$j=(J-1)/3>0$. Canonical $X=2^R$ with $R>3q$ gives $x=X-q>0$.
The source proves $w>q>F$, which makes $z_{\rm old}\ge2$; hence the
transport shift $z=z_{\rm old}-1$ is positive.

Supply $Y$ directly and set $r=\tau-\ell-s$. The first norm and
$\ell>k$ give $\tau<\ell+k$ and
\[
 \tau-\ell=\frac{\ell k+1}{\tau+\ell}
       >\frac{k}{3}>\frac{XY}{3}>s.
\]
Thus $r>0$ and $A_0=0$. The positive multiple $v=k-2n$ of $4XY^2$
satisfies $M\mid v$. Set
\[
 \gamma=\gamma_{\rm old}+v/M-w.
\]
This is positive for a particularly simple reason:
\[
 \frac vM\ge\frac{4XY^2}{4q^2Y}
       =\frac{wY}{q}\ge wq^2>w.
\]
The resulting $D_*$ equals $Mg$, proving $A_m=0$.

The auxiliary coordinates use the squared-modulus reconstruction already
proved in package~58. For $R=2r_0+1$ with $r_0$ odd, the elementary odd
Pell polynomials satisfy
\begin{gather*}
 \chi_t(2r+1)=tQ_r(t^2),\qquad Q_r(0)=(-1)^r(2r+1),\\
 Q_r(1-A^2)=(-1)^r\psi_A(2r+1).
\end{gather*}
Take $V_* =\chi_I(R)/I$ and $y_* =\psi_I(R)$. They are positive integers,
$V_*>1$, and $I^2V_*^2-(I^2-1)y_*^2=1$. The identities yield
$V_*\equiv-c\pmod{f^2}$ and $V_*\equiv-R\pmod c$.
Since $\gcd(c,f)=1$,
\[
 T=\frac{V_*+c+Rf^2}{cf^2}
\]
is a positive integer and gives the required $V=V_*$. Finally
$(I^2-1)(y_*^2-V_*^2)=V_*^2-1>0$, so $u=y_*-V_*>0$.
All five residuals vanish. This proves both existence implications in
Theorem~\ref{thm:pell26} without changing the fixed tuple.

\subsection{Exact degree and literal cost}

The cancellation identities used in the degree proof are
\begin{align*}
 A_0={}&q^6[r^2+\ell(2r-k)-1]+2q^3(\ell+r)Y+Y^2,\\
 D_{\rm old}={}&X+ac+\gamma H,\qquad L=-v+SX,\\
 A_m={}&M^2[(X+\gamma H)^2+2ac(X+\gamma H)-Hc^2-1]\\
       &\hspace{15mm}+2MD_{\rm old}HL+H^2L^2.
\end{align*}
These are polynomial-ring identities, not relations valid only on solutions.
They give the exact residual degrees
\[
 (\deg A_0,\deg A_m,\deg A_t,\deg A_a,\deg A_s)
       =(11,13,3,13,10).
\]
Put $q_0=3\beta j$, $X_0=x+q_0$, $k_0=\eta+\zeta$, and
$R_4=q_0^3(q_0-F)$. The two highest residual forms are
\begin{align*}
 L_m&=32q_0^2k_0X_0^2Y^4(4q_0^2Y\gamma+R_4),\\
 L_a&=2X_0^2Y^2f^4uR_4.
\end{align*}
Hence the degree-$26$ homogeneous form is $L_m^2+L_a^2$.
The coefficient of $x^4Y^4f^8u^2j^6F^2$ is exactly $4(3\beta)^6>0$;
$L_m^2$ contains no $f$ or $u$ and cannot change it. This establishes
the exact degree for the authentic numerical instance.

The evaluator has $49$ multiplications, $27$ additions and $16$
subtractions, totaling $92$ live binary operations. Scalar multiplication
is charged, every fixed power is expanded into multiplications, and all
fourteen inputs are live. This arithmetic metric is separate from bit
complexity and the height of the expanded coefficients.

\subsection{A degree-30 companion with 85 operations}

Retain the positive inputs
\[
 J,F,\alpha,z,f,i,T,w,Y,r,\eta,\zeta,u,\gamma,
\]
and put $q=\beta J+1$, $X=wq$. Use the same common expressions, but replace
the main carrier and transport by
\begin{align*}
 M&=4qY,\qquad v=k-R-1,\\
 D_*&=M(X+ac)+H(M\gamma-v),\\
 A_t&=(K_0+w)C-F-z(q-1).
\end{align*}
Keep $A_0,A_m,A_a,A_s$ and the five-square sum, calling it $P_{14,30}$.
The literal circuit is in \path{pell/circuits/p14_degree30.json}.

There is no restriction modulo $3$ on $J$ here. Once $q^3\mid Y$, every
prime common to $qY$ and $H$ is $3$. If $3\mid q$, then $27\mid Y$ and
$H\equiv3\pmod{27}$, so its $3$-adic valuation is exactly one. If
$3\nmid q$, then $\gcd(q,H)=1$ and $\gcd(Y,H)\mid3$. As $H$ is odd,
\[
 \gcd(4qY,H)\le3,
 \qquad\frac M{\gcd(M,H)}\ge\frac{4qY}{3}>R+1.
\]
Also $M\mid4XY^2$. The same ordered index proof gives $p=R$ and
$2n=R+1$, then $h=v/E>0$ and $\gamma_{\rm old}=\gamma-v/M>0$.
Completeness uses $\gamma=\gamma_{\rm old}+v/M>0$; both gap
constructions remain unchanged. Thus it has the same consistency
equivalence.

Its residual degrees are $(12,14,2,15,10)$. With
$q_0=\beta J$, $X_0=wq_0$ and $R_4=q_0^3(q_0-F)$, the highest form is
\[
 [P_{14,30}]_{30}=4X_0^4Y^4f^8u^2R_4^2.
\]
The coefficient of $w^4J^{10}Y^4f^8u^2F^2$ is $4\beta^{10}>0$.
The evaluator has $45$ multiplications, $24$ additions and $16$
subtractions, totaling $85$ operations.

\subsection{Verification and remaining formalization}

The independent checker uses the standard Python library and does not
import the circuit builders. It proves three local polynomial identities
by exact sparse algebra, checks all gate dependencies and counts, and
compares the DAGs with separately transcribed formulas at $200$ modular
assignments. It also checks $48$ exact polynomial rays modulo two primes
against the analytic leaders using the actual fixed tuple, exercises
$10{,}584$ bounded carrier cases, and rejects a transcription mutation
for each circuit. The ordinary mathematical degree proof above supplies
the universal upper bound and nonzero leading coefficient; the modular
rays are implementation checks. The same checks pass under Python's
\texttt{-O} mode.

The inherited coefficient audit independently rebuilds the four numerals,
verifies $66{,}062$ finite basis interactions and their provenance, tests
$38{,}663$ one-bit corruptions, and verifies $87{,}380$ small complement
models. These checks do not replace the unbounded native compiler proof.
An independent internal mathematical review of the new carrier is included
as \texttt{pell/INDEPENDENT\_PROOF\_AUDIT.md}; it is AI-assisted review,
not external human peer review.

The next formalization target is the two-stage carrier lemma with its exact
order of recovery: first the modulus $S$, then the index identity $p=R$,
then the full modulus $M$ and $q\mid X$. Potential further reductions must
preserve this order and justify all positive replacement witnesses.
Whether an auxiliary root can remove another coordinate without losing
the degree advantage remains open.
